\section{Introduction}\label{sec:intro}

People take turns quickly. In recordings of ten languages, answers to yes/no questions most often began
within about 200\,ms of the end of the question~\citep{stivers2009universals}, which implies that
listeners plan their response while the other person is still talking~\citep{levinson2015timing}.
Listeners also signal attention while the other person talks, give way when interrupted, and do not
treat every pause as an invitation to speak~\citep{sacks1974turntaking,skantze2021review}. A voice agent
that must \emph{do} something on a call, such as look up an order or change a booking, faces a harder
version of this problem: the work may take seconds or minutes, the caller keeps talking while it runs,
and some of what the caller says changes the work (``actually, make it eight'').

A classical cascade of speech-to-text, a language model and text-to-speech answers only after the whole
chain has run. On Full-Duplex-Bench v3, a benchmark of tool use under real disfluency, a
Whisper$\to$GPT-4o$\to$TTS cascade spoke its first words 8.78\,s after the end of the user's speech, later
than any other system tested, and completed the task after 10.12\,s~\citep{lin2026fdbv3}. End-to-end full-duplex models remove most of that delay but
trade it for floor control: on the Full-Duplex-Bench v1 leaderboard Moshi starts its response 0.265\,s
after the user stops but takes the floor in 98.5\% of synthetic pauses, while the GPT-Realtime entry does
so in 1.0\% of them and responds after 1.470\,s~\citep{fdb_repo,lin2025fdbv1}. Neither end of this
trade-off suits an agent that must also act.

In 2026 a third design became prominent. Thinking Machines Lab's \emph{interaction models} pair a
real-time interaction model, which processes audio in 200\,ms micro-turns, with an asynchronous background
model for reasoning and tool use that shares its context; the interaction model is a 276B-parameter
mixture-of-experts (12B active) trained from scratch, and the system was announced as a research
preview~\citep{tml2026interaction}. OpenAI's GPT-Live
delegates search and reasoning to a backend model behind an asynchronous
boundary~\citep{openai2026gptlive,openai2026gptlivesys}, Google's Gemini 3.8 Live runs tools in the
background while it keeps talking~\citep{google2026g38}, and open frameworks have added asynchronous tool
calls~\citep{livekit2026async,pipecat_fc}. The front/background split is therefore not new. What is
missing is an implementation that an operator can run on its own hardware, with its own choice of speech
vendors, models and agents, and inspect.

This paper describes \emph{interaction mode}, released in Bud WaaV Gateway 1.1.0 on 27 October
2026~\citep{waav_repo,waav_imdoc}. WaaV is an open-source (Apache-2.0) real-time voice gateway written in
Rust. Interaction mode adds a small \emph{front model} (Qwen3.5-2B~\citep{qwen35_2b} on
llama.cpp~\citep{llamacpp}) that names one explicit \emph{act} for every caller turn and acknowledges at
once; a \emph{task layer} that runs each request on any language model, an orchestrator or remote agents
and relays the answer verbatim; a \emph{duplex audio layer} that tracks the floor, yields to the caller and
ignores its own echo; and a \emph{policy layer} that guarantees that the words that bind (confirmations and
reports of what was done) come from templates and from the background's own text, never from the small
model. Our contributions are:
\begin{itemize}
  \item an open, self-hostable design in which a 2B front model makes an explicit, grammar-constrained,
    logged decision per caller turn while the gateway owns the binding words
    (Sections~\ref{sec:sys-front}--\ref{sec:sys-policy});
  \item background work as first-class tasks with revisions, committed-effect accounting, speculative
    start, an update inbox, follow-ups and life after the call, behind one open task protocol over
    webhooks, A2A, MCP and the OpenAI Responses API (Sections~\ref{sec:sys-tasks}--\ref{sec:sys-backgrounds});
  \item duplex behaviour inside a cascade, with a timeline of what the caller heard that decides when a
    spoken answer binds, and confirmations by voice, screen or keypad (Sections~\ref{sec:sys-audio}
    and~\ref{sec:sys-hitl});
  \item the same mode over 44 speech vendors, 11 speech-to-speech integrations, LiveKit rooms and SIP calls
    (Section~\ref{sec:sys-surfaces}), with measurements on Full-Duplex-Bench v1 and v3
    (Section~\ref{sec:eval}) and a case study (Section~\ref{sec:case}).
\end{itemize}
We do not claim that the design beats end-to-end models on every axis. A cascade has a latency floor
(the time to settle the end of a turn plus the time to the first synthesised audio) that a single model
trained on audio does not have. WaaV addresses \emph{perceived} latency, by acknowledging at once and
keeping the caller informed while work runs, and makes the agent's decisions explicit;
Section~\ref{sec:discussion} discusses the trade-off and the known limits of the release.
